\scanpage{76}
\mspage{2}
% archive foliation: 75

On a un candidat naturel (pour $\Sigma_{\deleted{J}}$ \inserted{$=\Sigma_J$}, \deleted{$D_J$} \inserted{$D=D_J$}, $t\in\Sigma^*$ donnés) comme droite $T$ munie de (5), savoir
\[
\deleted{(7)}\qquad T_0\simeq O_t(1)
\]
à cela près qu'il y a \uline{deux} isom. naturels
\[
\deleted{(8)}\qquad T_0^{\otimes2}=O_t(2)\xrightarrow{\sim}T_{\Sigma,t}\simeq O_t(2)(\underline\omega),
\]
correspondant aux deux éléments $\omega$ de $\underline\omega$, savoir
\[
(7)\qquad\deleted{\unsure{t}}\unsure{\xi|\zeta}\longmapsto\unsure{\xi|\zeta}\wedge\omega
\]
\margin{$T_0^{\otimes2}\xrightarrow{\sim}O_t(2)\simeq T_{\Sigma,t}\simeq\underline O_{\unsure{t}}(2)(\underline\omega)$}

On va lever cette ambiguïté, en remplaçant $T_0=O_t(1)$ par une variante tordue. La tordre par un torseur sous $\mu_2$ \inserted{$=\mu_2(k)$} ne nous avance pas, car \deleted{$\unsure{O}$} $T_0^{\otimes2}$ n'aura pas changé (à isom can. près). On va la tordre par un torseur (qu'on précisera) $Q_0$ sous $\mu_4=\mu_4(k)$ (qui peut donc s'interpréter comme l'ens des sommets d'un \uline{carré}, dont $\mu_4$ est le groupe des rotations, avec $\{i,-i\}=\mu_4^*$ comme \uuline{paire} de génér. privilégiés\ldots).

On posera donc
\[
(8)\qquad T_0=\underline O_t(1)\wedge_{\mu_4}Q_0
\]
d'où
\[
(9)\qquad T_0^{\otimes2}\simeq\underline O_t(2)\wedge_{\mu_2}D_{Q_0}
\]
où
\[
(10)\qquad D_{Q_0}=Q_0/\{\pm1\}
\]
est l'ens des deux diagonales du carré $Q_0$. Supposons donc qu'on se donne un carré \inserted{\unsure{combinatoire} (``sommets'')} $Q_{0S}$ avec des isom. \deleted{\unsure{$\operatorname{Aut}Q_0$} est le groupe des}
\[
(11)\qquad\left\{
\begin{array}{lll}
\underline\omega_{Q_0}&\text{(ens des orientations de $Q_0$)}&\simeq\mu_4^*=\{i,-i\}\\
D_{Q_0}&\text{(ens des diagonales de $Q_0$)}&\simeq\underline\omega
\end{array}\right.
\]
